\documentclass[10pt,a4paper]{amsart}
\usepackage[margin=19mm]{geometry}
\usepackage{amsmath,amssymb,mathtools}
\usepackage[scr=rsfs,cal=euler]{mathalfa}
\usepackage{xfrac}
\usepackage[unicode,hidelinks,bookmarks=false]{hyperref}
\newcommand{\ie}{i.e.\ }
\newcommand{\cf}{cf.\ }
\newcommand{\resp}{respectively\ }
\newcommand{\Subsection}{\subsection}
\providecommand{\nobreakdash}{\nobreak\hbox{-}}
\setlength{\emergencystretch}{2em}
\multlinegap0pt
\usepackage{amssymb}
\usepackage{mathtools}
\usepackage{bm}
\usepackage[shortlabels]{enumitem}
\usepackage{xcolor}
\usepackage{tikz}
\newtheorem{theorem}{Theorem}
\newcommand{\ml}[2]{\mathcal{L}(#1\!\to\!#2)}
\title{Minimax Quantile Bounds via Information Measures}
\author{Amedeo Roberto Esposito}
\date{Source: arXiv:2608.20857v2 (CC BY 4.0)}
\begin{document}
\maketitle

\noindent\emph{Source and licence.} Minimax Quantile Bounds via Information Measures. Version arXiv:2608.20857v2 (CC BY 4.0), distributed by arXiv under the Creative Commons Attribution 4.0 International licence, http://creativecommons.org/licenses/by/4.0/, as recorded in the arXiv licence metadata for this exact version and shown on the abstract page https://arxiv.org/abs/2608.20857v2. Full licence notice: \texttt{licenses/arxiv-2608.20857-CC-BY-4.0.txt}.\par\smallskip

\noindent\textit{Editorial scope.} Counted unit: the article's theorem thm:gwsbmFiniteLeakage (source0.tex lines 1625--1689). The article's complete original proof of it is delivered with it (source0.tex lines 1691--2029, one \texttt{proof} environment of 339 lines). No mathematics is written, repaired or re-derived: the statement and the proof are the article's own lines, in the article's own notation, and no retained line is substituted or re-flowed.  No further source-local context is needed: every cross-reference of every retained line resolves inside the two delivered spans.  Nothing was omitted from any retained span, so the package carries no omission record.  No retained line contains a citation, so the wrapper carries no bibliography and no bibliography entry.  No retained line contains a figure environment, an image-inclusion command or drawing code, so no image is delivered and the package declares no asset dependency.  Editorial formatting only, in addition: an AMS \texttt{amsart} preamble that restates the article's own declarations and macros used by the retained lines, namely the environments \texttt{theorem} on separate counters, together with the standard packages those lines need.  The retained environments print in this document as Theorem 1 in order of appearance, whereas the article prints the same items with the numbering of its own full document.  The complete native source of the article as distributed in the arXiv e-print for this exact version is bundled unchanged as \texttt{sources/208206/source0.tex}.
\begin{theorem}
\label{thm:gwsbmFiniteLeakage}
For every integer \(1\leq L\leq\lfloor m/2\rfloor\) and every \(t>0\),
one has
\begin{align}
&
1-
P\left(
U_m+V_m-\frac{2G}{\sqrt{n-2}}
\geq
2\vartheta_n-\varepsilon_{n,L}(t)
\right)
-r_n(t)-B_{n,L}
\notag\\
&\leq
\frac{\exp\{\ml{W}{A}\}}{|\Theta_n|}
\notag\\
&\leq
1-
P\left(
U_m+V_m-\frac{2G}{\sqrt{n-2}}
\geq
2\vartheta_n+\varepsilon_{n,1}(t)
\right)
+r_n(t),
\label{eq:gwsbmFiniteLeakageSandwich}
\end{align}
Here
\[
\begin{aligned}
r_n(t)&:=n(n-1)\overline\Phi(t),\\
B_{n,L}&:=
\sum_{k=L+1}^{\lfloor m/2\rfloor}
\binom{m}{k}^{2}
\overline\Phi\!\left(2\lambda_n\sqrt{k(m-k)}\right),\\
\vartheta_n&:=\frac{(n-1)\lambda_n}{\sqrt{n-2}},
&
\varepsilon_{n,L}(t)&:=
\frac{2(2L-1)(\lambda_n+t)}{\sqrt{n-2}}.
\end{aligned}
\]

Consequently,
\begin{align}
&
P\left(
U_m+V_m-\frac{2G}{\sqrt{n-2}}
\geq
2\vartheta_n+\varepsilon_{n,1}(t)
\right)
-r_n(t)
\notag\\
&\leq
\mathcal R_n^\star
\notag\\
&\leq
P\left(
U_m+V_m-\frac{2G}{\sqrt{n-2}}
\geq
2\vartheta_n-\varepsilon_{n,L}(t)
\right)
+r_n(t)+B_{n,L}.
\label{eq:gwsbmFiniteRiskSandwich}
\end{align}
\end{theorem}

\begin{proof}
For every $[\sigma]\in\Theta_n$, let $p_\sigma$ denote the density of
$P_\sigma$ with respect to Lebesgue measure on
$\mathbb R^{\binom{n}{2}}$. Since $W$ is uniformly distributed over
$\Theta_n$, the exponential of Maximal Leakage is
\begin{equation}
\exp\{\ml{W}{A}\}
=
\int_{\mathbb R^{\binom{n}{2}}}
\max_{[\sigma]\in\Theta_n}
p_\sigma(a)
\,\mathrm da.
\label{eq:gwsbmLeakageIntegral}
\end{equation}

For every $[\sigma]\in\Theta_n$, define its maximum-likelihood region by
\begin{equation}
\mathcal D_\sigma
:=
\left\{
a:
p_\sigma(a)>
p_{\sigma'}(a)
\text{ for every }
[\sigma']\neq[\sigma]
\right\}.
\label{eq:gwsbmMLRegion}
\end{equation}
Likelihood ties have Lebesgue measure zero, and therefore
\begin{align}
\exp\{\ml{W}{A}\}
&=
\sum_{[\sigma]\in\Theta_n}
\int_{\mathcal D_\sigma}
p_\sigma(a)
\,\mathrm da.
\label{eq:gwsbmLeakageDecisionRegions}
\end{align}
The action of the permutation group on $\Theta_n$ is transitive, and
simultaneously permuting the rows and columns of $A$ maps
$\mathcal D_\sigma$ onto the corresponding maximum-likelihood region.
Consequently, all the integrals in
\eqref{eq:gwsbmLeakageDecisionRegions} are equal. For the fixed
$[\sigma^\star]\in\Theta_n$,
\begin{equation}
\frac{\exp\{\ml{W}{A}\}}{|\Theta_n|}
=
\int_{\mathcal D_{\sigma^\star}}
p_{\sigma^\star}(a)
\,\mathrm da
=
P_{\sigma^\star}
\left(
A\in\mathcal D_{\sigma^\star}
\right).
\label{eq:gwsbmNormalisedLeakageRegion}
\end{equation}
Thus, the normalised exponential of Maximal Leakage is the
$P_{\sigma^\star}$-measure of the region on which the true labelling
maximises the likelihood.

We now characterise this region through balanced exchanges. Let
\(U\subseteq C_+\) and \(V\subseteq C_-\) satisfy \(|U|=|V|=k\),
and let $\sigma^{U,V}$ be the labelling obtained from
$\sigma^\star$ by reversing the signs of the vertices in $U\cup V$.
Define
\begin{equation}
T_{U,V}
:=
\sum_{\substack{1\leq i<j\leq n\\
|\{i,j\}\cap(U\cup V)|=1}}
Y_{ij}.
\label{eq:gwsbmExchangeStatistic}
\end{equation}
A direct calculation of the Gaussian likelihood ratio gives
\begin{equation}
\log
\frac{
p_{\sigma^{U,V}}(A)
}{
p_{\sigma^\star}(A)
}
=
-2\lambda_n T_{U,V}.
\label{eq:gwsbmExchangeLikelihoodRatio}
\end{equation}
Since $\lambda_n>0$, the competing labelling has likelihood at least
as large as the true labelling if and only if
\begin{equation}
T_{U,V}\leq0.
\label{eq:gwsbmDetrimentalExchange}
\end{equation}
Every element of $\Theta_n$ can be represented in this way, and,
because a labelling and its global sign reversal are identified, it is
enough to consider \(1\leq k\leq\lfloor m/2\rfloor\).
Up to an event of probability zero,
\begin{equation}
\mathcal D_{\sigma^\star}
=
\bigcap_{k=1}^{\lfloor m/2\rfloor}
\bigcap_{\substack{U\subseteq C_+,\,V\subseteq C_-\\
|U|=|V|=k}}
\left\{
T_{U,V}>0
\right\}.
\label{eq:gwsbmMLRegionIntersections}
\end{equation}

Let
\begin{equation}
\mathcal E_t
:=
\left\{
\max_{1\leq i<j\leq n}|Z_{ij}|\leq t
\right\}.
\label{eq:gwsbmBoundedNoiseEvent}
\end{equation}
A union bound gives
\begin{align}
P_{\sigma^\star}
\left(
\mathcal E_t^c
\right)
&\leq
2\binom{n}{2}\overline\Phi(t)\\
&=
n(n-1)\overline\Phi(t)\\
&=
r_n(t).
\label{eq:gwsbmBoundedNoiseProbability}
\end{align}

We first derive an upper bound on
$\exp\{\ml{W}{A}\}/|\Theta_n|$. Let
$i^\star\in C_+$ and $j^\star\in C_-$ satisfy
\(-R_{i^\star}=M_{+,n}\) and \(-R_{j^\star}=M_{-,n}\).
For the exchange of $i^\star$ and $j^\star$,
\begin{align}
T_{\{i^\star\},\{j^\star\}}
&=
D_{i^\star}+D_{j^\star}-2Y_{i^\star j^\star}\\
&=
\sqrt{n-2}
\left(
2\vartheta_n-M_{+,n}-M_{-,n}
\right)
-
2\left(
\lambda_n+Z_{i^\star j^\star}
\right).
\label{eq:gwsbmMostDetrimentalOneExchange}
\end{align}
If \(M_{+,n}+M_{-,n}\geq
2\vartheta_n+\varepsilon_{n,1}(t)\) and $\mathcal E_t$ occurs, then
\begin{align}
T_{\{i^\star\},\{j^\star\}}
&\leq
-2(\lambda_n+t)
-
2(\lambda_n-t)\\
&=
-4\lambda_n
<
0.
\end{align}
It follows from~\eqref{eq:gwsbmMLRegionIntersections} that
\begin{equation}
\left\{
M_{+,n}+M_{-,n}
\geq
2\vartheta_n+\varepsilon_{n,1}(t)
\right\}
\cap
\mathcal E_t
\subseteq
\mathcal D_{\sigma^\star}^c.
\label{eq:gwsbmLeakageUpperInclusion}
\end{equation}
Using~\eqref{eq:gwsbmNormalisedLeakageRegion} and
\eqref{eq:gwsbmBoundedNoiseProbability},
\begin{align}
\frac{\exp\{\ml{W}{A}\}}{|\Theta_n|}
&=
P_{\sigma^\star}
\left(
A\in\mathcal D_{\sigma^\star}
\right)\\
&\leq
1-
P_{\sigma^\star}
\left(
M_{+,n}+M_{-,n}
\geq
2\vartheta_n+\varepsilon_{n,1}(t)
\right)
+r_n(t).
\label{eq:gwsbmLeakageUpperBound}
\end{align}

We next derive the lower bound. Suppose that
$T_{U,V}\leq0$ for some \(U\subseteq C_+\) and \(V\subseteq C_-\)
with \(|U|=|V|=k\leq L\).
Writing $U\cup V$ for the set of $2k$ exchanged vertices, one has
\begin{align}
\sum_{i\in U\cup V}D_i
&=
T_{U,V}
+
2
\sum_{\substack{1\leq i<j\leq n\\
i,j\in U\cup V}}
Y_{ij}.
\label{eq:gwsbmVertexScoreDecomposition}
\end{align}
On $\mathcal E_t$, the second sum contains
\(\binom{2k}{2}=k(2k-1)\) terms, each bounded above by
$\lambda_n+t$. Since $T_{U,V}\leq0$, it follows that
\(\sum_{i\in U\cup V}D_i\leq2k(2k-1)(\lambda_n+t)\).
Using \(D_i=(n-1)\lambda_n+\sqrt{n-2}\,R_i\),
we obtain
\begin{equation}
-\sum_{i\in U\cup V}R_i
\geq
2k\vartheta_n
-
\frac{
2k(2k-1)(\lambda_n+t)
}{
\sqrt{n-2}
}.
\label{eq:gwsbmDetrimentalSmallExchangeLower}
\end{equation}
Moreover,
\(-\sum_{i\in U\cup V}R_i\leq kM_{+,n}+kM_{-,n}\).
Consequently, every detrimental exchange with $k\leq L$, on
$\mathcal E_t$, implies
\begin{align}
M_{+,n}+M_{-,n}
&\geq
2\vartheta_n
-
\frac{
2(2k-1)(\lambda_n+t)
}{
\sqrt{n-2}
}\\
&\geq
2\vartheta_n-\varepsilon_{n,L}(t).
\label{eq:gwsbmSmallExchangeImplication}
\end{align}

It remains to control exchanges with $k>L$. For fixed
$U\subseteq C_+$ and $V\subseteq C_-$ satisfying
$|U|=|V|=k$, the statistic $T_{U,V}$ is the sum of
$4k(m-k)$ independent random variables with distribution
$\mathcal N(\lambda_n,1)$. Hence
\(P_{\sigma^\star}(T_{U,V}\leq0)
=\overline\Phi(2\lambda_n\sqrt{k(m-k)})\).
There are $\binom{m}{k}^2$ possible pairs $(U,V)$. Therefore,
\begin{align}
&P_{\sigma^\star}
\left(
T_{U,V}\leq0
\text{ for some }
U,V
\text{ with }
L<|U|=|V|
\leq
\left\lfloor\frac m2\right\rfloor
\right)\\
&\qquad\leq
\sum_{k=L+1}^{\lfloor m/2\rfloor}
\binom{m}{k}^2
\overline\Phi
\left(
2\lambda_n\sqrt{k(m-k)}
\right)\\
&\qquad=
B_{n,L}.
\label{eq:gwsbmLargeExchangeProbability}
\end{align}

Combining~\eqref{eq:gwsbmMLRegionIntersections},
\eqref{eq:gwsbmSmallExchangeImplication}, and
\eqref{eq:gwsbmLargeExchangeProbability}, one obtains
\begin{align}
\mathcal D_{\sigma^\star}^c
\subseteq{}&
\left\{
M_{+,n}+M_{-,n}
\geq
2\vartheta_n-\varepsilon_{n,L}(t)
\right\}
\cup
\mathcal E_t^c
\notag\\
&\cup
\left\{
T_{U,V}\leq0
\text{ for some }
U,V
\text{ with }
L<|U|=|V|
\leq
\left\lfloor\frac m2\right\rfloor
\right\}.
\end{align}
Using~\eqref{eq:gwsbmNormalisedLeakageRegion},
\eqref{eq:gwsbmBoundedNoiseProbability}, and
\eqref{eq:gwsbmLargeExchangeProbability}, it follows that
\begin{align}
\frac{\exp\{\ml{W}{A}\}}{|\Theta_n|}
&\geq
1-
P_{\sigma^\star}
\left(
M_{+,n}+M_{-,n}
\geq
2\vartheta_n-\varepsilon_{n,L}(t)
\right)
\notag\\
&\qquad
-r_n(t)-B_{n,L}.
\label{eq:gwsbmLeakageLowerBound}
\end{align}

Finally, using
\(M_{+,n}+M_{-,n}\overset{\mathrm d}{=}
U_m+V_m-2G/\sqrt{n-2}\)
in~\eqref{eq:gwsbmLeakageUpperBound} and
\eqref{eq:gwsbmLeakageLowerBound} proves
\eqref{eq:gwsbmFiniteLeakageSandwich}.

Since the maximum-likelihood estimator is an equaliser and minimax,
\(\mathcal R_n^\star
=1-\exp\{\ml{W}{A}\}/|\Theta_n|\).
Taking complements in the Maximal Leakage bounds gives
\eqref{eq:gwsbmFiniteRiskSandwich}.
\end{proof}

\end{document}
